\chapter{Lezioni Apprese e Riflessioni Conclusive}

In questo capitolo conclusivo vengono affrontate tre domande chiave per riflettere sui risultati ottenuti e sull'impatto delle pratiche di sostenibilità applicate durante lo sviluppo e il refactoring del progetto 2Chance.

\section*{Cosa hai imparato sull'ingegneria del software sostenibile?}
Durante questo progetto, ho capito che rendere il software sostenibile non è un'attività isolata, ma è un requisito non funzionale che deve essere tenuto in considerazione durante ogni fase del ciclo di vita del software. Ho imparato che l'efficienza energetica del codice è strettamente legata alle scelte architetturali, come l'ottimizzazione delle query al database e l'eliminazione di anti-pattern strutturali (ad esempio le invocazioni ripetute nei cicli o il "Problema N+1"). Inoltre, ho acquisito consapevolezza sull'importanza della \textit{Software Supply Chain} e della sostenibilità legale: utilizzare strumenti per la generazione della SBOM e per il monitoraggio delle licenze (come FOSSA) è fondamentale per garantire che un prodotto sia etico, trasparente e privo di rischi legali. Infine, la sostenibilità sociale evidenzia quanto siano essenziali le dinamiche all'interno del team, la comunicazione e un sano approccio alla condivisione per prevenire i \textit{community smells}.

\section*{Quali trade-off hai affrontato (es. performance vs. manutenibilità)?}
L'adozione di pratiche di sostenibilità ha introdotto alcune sfide. In primo luogo, è emerso un chiaro trade-off tra \textbf{performance e manutenibilità}: per garantire una maggiore efficienza energetica e risolvere i colli di bottiglia nei \texttt{Model DAO}, è stato necessario estrarre le logiche SQL dai cicli per raggrupparle in transazioni \textit{batch}. Sebbene questo intervento abbia risolto il problema delle chiamate N+1 abbattendo radicalmente i consumi, il codice risultante è diventato più verboso e complesso da mantenere rispetto all'utilizzo di metodi ingenui e diretti. Un secondo compromesso fondamentale ha riguardato l'\textbf{investimento iniziale} di tempo. Configurare rigorosamente i benchmark hardware tramite strumenti come EnergiBridge e impostare delle solide pipeline di CI/CD rallenta lo sviluppo nelle sue fasi iniziali; tuttavia, si tratta di una scelta indispensabile, poiché garantisce un enorme guadagno architetturale e una maggiore scalabilità nel medio-lungo termine.

\section*{In che modo questo progetto si collega alla qualità del software a lungo termine?}
L'applicazione di pratiche sostenibili per il progetto 2Chance porterà miglioramenti a lungo termine poiché risolvendo alla radice le violazioni di eco-design, il sistema è diventato nativamente più scalabile: consumando meno risorse (CPU, RAM e I/O di rete) per singola operazione, l'applicazione potrà supportare un carico di utenti maggiore senza richiedere aggiornamenti hardware precoci, riducendo così i futuri e-waste.
Parallelamente, l'integrazione di unit e integration testing automatizzati, l'analisi statica della sostenibilità con Creedengo e controlli di licenza normativi con FOSSA direttamente nella pipeline automatizzata assicura che il codice non vada incontro ad un decadimento della sostenibilità ambientale, tecnica e legale.
